\section{Background}
\label{sec:background}

\subsection{Vector databases and k-nearest neighbor search}

A vector database stores high dimensional vectors that represent
unstructured records such as text passages, code snippets, or images
embedded by a learned model. The dominant query primitive is
$k$-nearest neighbor search: given a query vector, the system returns
the $k$ corpus vectors closest to it under a fixed distance metric, for
example cosine similarity or Euclidean distance. Because exhaustively
scanning every vector is prohibitive at corpus sizes of tens or
hundreds of millions, modern systems rely on approximate nearest
neighbor indices that trade a small recall loss for one or two orders
of magnitude in latency.

The literature has produced several families of approximate nearest
neighbor index, each making a different trade-off between build
time, memory footprint, query latency, and recall.
\emph{Hashing based} indices, exemplified by locality sensitive
hashing, bucket vectors so that similar points collide with high
probability, trading recall for very fast lookups.
\emph{Inverted file} indices cluster the corpus into Voronoi cells
and probe only the cells closest to the query at search time, and
are often combined with product quantization to compress each
vector down to a few bytes and shrink the resident memory footprint
at the cost of some recall. \emph{Graph based} indices, of which
HNSW~\cite{malkov2018hnsw} is the most widely deployed example,
treat each vector as a node linked to its approximate neighbors and
answer queries by greedy walks on the graph. Graph indices
currently dominate published benchmarks for in-memory approximate
nearest neighbor search. Across all of these families the
index ultimately determines which vectors a query inspects, and
therefore determines the query performance.

\subsection{LSM based search architecture}

Vector databases that ingest data continuously cannot rebuild a single
monolithic index on every write. Systems such as
Milvus~\cite{wang2021milvus} and Manu~\cite{guo2022manu} therefore
adopt a log structured merge (LSM) style layout. Incoming vectors are
first buffered in a mutable \emph{growing segment} held in memory.
When the growing segment reaches a target size it is flushed to object
storage as an immutable \emph{sealed segment}, after which a per
segment vector index, typically an HNSW graph, is built over the
flushed vectors. A collection therefore exists on disk as a sequence of
sealed segments, each carrying its own data file and its own index
structure, plus at most one growing segment in memory.

The query path mirrors this layout. A search request is fanned out to
every segment in the collection, including the growing one. Each
segment independently searches its local index and returns its own
top $k$ candidates, and a coordinator merges the per segment results
into the global top $k$ that is returned to the client. Because every
segment is consulted on every query, the per query cost scales with
the number of segments and with the size of each segment's index,
which means the segment layout, how many segments exist, how large
they are, and which vectors they contain, directly determines query
latency and the working set held in memory or page cache.
%Two further mechanisms reshape this layout in practice. First, Milvus jitters the
% flush threshold so that successive sealed segments do not all reach
% the seal point at exactly the same vector count, which prevents
% synchronized flushes across data nodes. Second, a background
% compactor periodically merges small or partially deleted sealed
% segments into larger ones, rewriting their files in the process.
% Both mechanisms move vectors across segment boundaries after ingestion
% and are central to the storage behavior we examine in the remainder
% of this report.
